\label{app:closureappendix}

Every paper in this field defends a low precision figure the same way: the references are
incomplete, so a chemically real product that no curator recorded is counted as an error. This
paper has leaned on that defence too (\S\ref{sec:tame}). It is untestable in general, because what
was never annotated is not observable. One component of it is testable, and this appendix measures
it.

Metabolism composes. The corpus records conversions as edges, substrate to product, and those edges
form a graph rather than a list of independent rows. If it records $A \to B$ and, elsewhere,
$B \to C$, then it asserts a two-step route from $A$ to $C$ without any curator having written the
edge $A \to C$. A method that emits $C$ for $A$ is scored wrong, and whether that is right depends
on a fact about how the corpus was assembled rather than on the chemistry.

\paragraph{The corpus is a graph and it is not transitively closed.} Node identity is the
tautomer-aware InChIKey used throughout, so the structure is not an artifact of how a curator drew a
molecule, and the graph is taken over all three clean splits, because the corpus is what asserts the
second edge while a split only asks about the first. It carries $21{,}964$ edges and $7{,}296$
two-step compositions, which resolve to $7{,}234$ distinct composed pairs. The direct edge is
annotated for $302$ of them. The remaining $6{,}932$ are routes the corpus asserts in two steps and
denies in one.

\paragraph{What that is worth to a method, measured.} If the defence carried weight, a method's
scored-wrong output would be full of these. It is not. Of everything the metric counts wrong, the
share the corpus itself reaches from that substrate in two steps is $0.0016$ $[0.0007,0.0028]$ for
GRAIL, $0.0035$ $[0.0021,0.0052]$ for MetaPredictor and $0.0015$ $[0.0009,0.0023]$ for SyGMa. The
quantity is method-dependent in the way the argument predicts, and two of the three pairwise
differences have intervals excluding zero, the largest being MetaPredictor over SyGMa at $+0.0020$
$[+0.0008,+0.0034]$. It is also two orders of magnitude too small to matter. Crediting a method for
everything the annotation graph reaches in two steps raises precision from $0.1060$ to $0.1075$ for
GRAIL, from $0.1061$ to $0.1090$ for MetaPredictor and from $0.0777$ to $0.0791$ for SyGMa; a third
step adds a further $0.0002$; and the ordering by F1 is the same at every depth.

\paragraph{Why it is worth reporting a null.} The point is not that annotation is complete; it is
that the one part of the incompleteness the data can audit is small, so a precision figure on this
corpus cannot be explained away by it. That removes an excuse rather than supplying one, including
from this paper, which has used it. It also settles a design question cheaply: a benchmark rebuilt
by closing its annotation graph under composition would turn $302$ annotated compositions into
$7{,}234$ and would not change what any of these methods is measured to do, because none of them
emits the difference.

\paragraph{The model's own output is a graph too, and it composes.} The corpus is not the only thing
here with edges. A method that predicts $A \to B$ and, separately, $B \to C$ has predicted a two-step
route from $A$ to $C$ without emitting $C$ for $A$, and a one-step evaluation scores that as a miss.
Composition is restricted to intermediates that are themselves substrates of this split, the only
molecules every method has a prediction to compose with, which leaves $389$ substrates for GRAIL,
$428$ for MetaPredictor and $367$ for SyGMa.

Composing enlarges the emitted set, and a larger set matches more for no good reason, so the
recovery is reported against a control that enlarges it by the same amount: composing through the
same number of substrates drawn at random. That control recovers $0.0000$ for all three methods, to
four decimals, which is what makes the rest a statement about the method rather than about set size.
Of the references a one-step evaluation records as missed, composing the method's own predictions
recovers $0.0086$ $[0.0041,0.0138]$ for GRAIL, $0.0243$ $[0.0169,0.0328]$ for MetaPredictor and
$0.0136$ $[0.0080,0.0202]$ for SyGMa, every one of them separated from its control. The quantity
differs by method: MetaPredictor recovers $0.0158$ $[0.0069,0.0249]$ more than GRAIL and $0.0107$
$[0.0028,0.0188]$ more than SyGMa, both separated.

It is not a policy. Recovering one reference costs $86$ additional candidates for MetaPredictor and
$187$ for GRAIL, so emitting the composition would destroy precision many times over what it buys.
What it is is a property of a method that recall does not report: how nearly its one-step output is
closed under its own dynamics. Two methods with the same recall can differ in it, and here they do.

\paragraph{One further reading.} The corpus is closed at $4.2\%$ and the test split alone at
$9.7\%$, on $390$ composed pairs against the corpus's $7{,}234$. The gap between the two is not
noise about depth: within one split, chains tend to be annotated together because they come from one
record, while a composition that crosses splits joins records from different sources that no curator
ever linked. Which of the two a substrate falls under is another undeclared property of the
population it sits in (\S\ref{sec:robust}).
